Small language models are attractive engines for local, private agents, but agents increasingly expose many tools at once.
We measure how four off-the-shelf 1.5--3B instruct models (Qwen2.5-1.5B/3B, Llama-3.2-3B and Qwen3.5-2B), quantized and served with Ollama on a 4\,GB laptop GPU, cope as the tool catalog grows from 1 to 64.
We use only existing benchmarks: BFCL catalogs built from BFCL's own tools, BFCL's native user queries and executable multi-turn environments, and ToolRet.
Accuracy is stable up to about four tools and then falls by up to \DropQwenThree{} points at 64 tools.
A failure taxonomy derived from the raw generated tokens shows that the loss comes mostly from \emph{not calling} and from \emph{invalid call syntax}, not from choosing the wrong tool.
Two measurement pitfalls matter as much as the model: the serving stack silently discards malformed calls and calls to unknown tools, and strict scoring hides a \LenDropLlama-point decline for a model that writes numbers as strings.
Showing the model only the top-3 tools retrieved by a 0.6B embedder recovers \RetGainQwenThreeFive--\RetGainQwenOneFive{} points.
On native data, wrong-tool errors stay rare, but one model calls a tool in \NatAbstInvLlama\% of the cases where none applies. In executable multi-turn episodes, no model completes more than a quarter of the tasks.
The full study runs in about four GPU-hours, and we release all configurations, raw outputs, analysis code and a dated lab notebook.
